\chapter{Adverse Selection and Toxicity}\label{ch:hf:adverse-selection-and-toxicity}

A market maker's fills from one class of trader lose almost a tick a share ten seconds later, and its fills from the others gain half a tick. Its
quotes are the same for both until it decides they should not be. In this chapter's simulated market 32\% of a one-lot market maker's fills come
from informed traders, and those fills mark out at $-0.95$ ticks a share against $+0.55$ for the rest. Four features of the book, visible while the
quote rests, forecast a fill's mark-out with a correlation of 0.18 out of sample; withdrawing the side the forecast marks as toxic cuts the informed
share of fills from 29\% to 18\%. How often to withdraw matters as much as whether to: at the natural threshold the market maker churns its orders
to the back of their queues.

\section{Mark-outs by counterparty, venue and order type}

Adverse selection (One Quant Book 1, chapter 1) is measured by mark-outs (One Quant Book 2, chapter 15; One Quant Book 7, chapter 23): the move of
a reference price after a fill, signed against the market maker. What makes the measure useful is grouping. A dealer answering requests for quote
knows its counterparty and can mark out each client (One Quant Book 9, chapter 24, prices clients this way). On an anonymous exchange the market
maker does not know who hit it, but it knows the venue, the order type that hit it (an immediate-or-cancel order routed across venues, a market
order, an order that swept several levels), the size, and the time of day. Battalio, Corwin and Jennings found, in proprietary and public US equity
data, that limit orders sent to venues paying large rebates were executed worse on several measures of quality, adverse selection among them: the
venue is a group worth marking out by.

In a simulation the counterparty's class is known. The simulator records, for each of the market maker's fills, whether the aggressor was an
informed trader, and the market maker never sees it; \cref{fig:hf:adverse-selection-and-toxicity:classes} splits a one-lot quoter's fills over six
twenty-minute sessions. Informed fills turn against the quoter within two seconds and keep going: $-0.95$ ticks a share at ten seconds (standard error
0.14). Uninformed fills keep their half-spread: $+0.55$ (0.10). The informed traders are 32\% of the fills and set the sign of the total, $+0.07$.

\begin{omfigure}
\begin{tikzpicture}
\begin{axis}[width=11cm,height=5.2cm,xlabel={\small seconds after the fill},ylabel={\small mark-out (ticks a share)},xmode=log,log basis x=10,
  xmin=0.4,xmax=70,xtick={0.5,1,2,5,10,20,60},xticklabels={0.5,1,2,5,10,20,60},ymin=-1.6,ymax=0.8,tick label style={font=\footnotesize},
  grid=major,grid style={gray!25},table/col sep=comma,legend style={font=\scriptsize,at={(0.5,-0.3)},anchor=north},legend columns=2]
\addplot[omThm,thick,mark=*,mark size=1.2pt] table[x=h,y=informed]{figdata/hft/06-adverse-selection-and-toxicity/classes.csv};
\addplot[omDef,thick,dashed,mark=square*,mark size=1.2pt] table[x=h,y=uninformed]{figdata/hft/06-adverse-selection-and-toxicity/classes.csv};
\draw[gray] (axis cs:0.4,0) -- (axis cs:70,0);
\legend{informed aggressor,uninformed aggressor}
\end{axis}
\end{tikzpicture}

\omcaption{Mark-out of a one-lot market maker's fills against the mid, by the class of the trader who hit it (the simulator's truth), over six
simulated sessions of twenty minutes; quantity-weighted. Data: \texttt{fig\_toxicity.py} on \texttt{hf\_toxicity.run}.}\label{fig:hf:adverse-selection-and-toxicity:classes}
\end{omfigure}

\section{Detecting informed flow}

The market maker cannot see the class of the next aggressor; it can see the circumstances in which informed traders arrive. Flow toxicity (One
Quant Book 7, chapter 9) is the name for the share of the flow that adversely selects liquidity providers. Easley, López de Prado and O'Hara built
VPIN, an estimate of it from order imbalance in volume time, and proposed it as a warning of toxicity-induced volatility. Andersen and Bondarenko
found VPIN a poor volatility predictor, highest only after the flash crash, with its predictive content coming from a mechanical relation with
trading intensity, and warned against adopting any stress metric before comparing it with simple benchmarks. The lesson for a market maker is to
predict what it cares about, the mark-out of its own next fill, and to judge the predictor by that.

In a large-tick book the circumstances are few. Chapter 5 showed that a side whose queue is small against the other's is about to be eaten; a
burst of aggressive volume against a side announces the same thing; a mid that has not moved for a while is more likely to be stale; a spread wider
than a tick says the book is being repriced. The chapter's features are these four, computed from the order-by-order feed as the quote rests
(\cref{lst:hf:adverse-selection-and-toxicity:features}): the aggressive volume against the side over the last two seconds (lots), the other side's
share of the two best queues, the logarithm of one plus the seconds since the mid changed, and whether the spread exceeds a tick.

\section{Toxicity scores}

\begin{definition}[Toxicity score]\label{def:hf:adverse-selection-and-toxicity:score}
A \emph{toxicity score}\index{toxicity score} is a market maker's real-time estimate of the adverse selection it would suffer if one side of its
quotes were filled now: here, the negative of a forecast of the fill's mark-out at a fixed horizon, from features observable while the quote rests,
fitted on the market maker's own past fills.
\end{definition}

Fitted by least squares on the 919 fills of six calibration sessions, the ten-second mark-out forecast is
\[
\widehat{\mathrm{mo}}=1.43-0.029\,\text{flow}-1.79\,\text{queue share of the other side}-0.072\ln(1+\text{quiet})+0.028\,\text{wide},
\]
in ticks a share. The other side's share of the queues carries almost all of it, as chapter 5 predicted. The forecast correlates 0.19 with the
realised mark-outs it was fitted on and 0.18 with those of six other sessions: small, and stable. Sorted into fifths by the forecast, the other
sessions' fills realise mark-outs from $-0.86$ to $+0.83$ ticks a share (\cref{fig:hf:adverse-selection-and-toxicity:calibration}); a correlation of
0.18 on fills whose mark-outs vary by several ticks is enough to separate them.

\begin{omfigure}
\begin{tikzpicture}
\begin{axis}[width=11cm,height=5.2cm,xlabel={\small forecast mark-out, fifth of fills},ylabel={\small ticks a share},xtick={1,2,3,4,5},
  xmin=0.5,xmax=5.5,ymin=-1.3,ymax=1.3,tick label style={font=\footnotesize},grid=major,grid style={gray!25},table/col sep=comma,
  legend style={font=\scriptsize,at={(0.5,-0.3)},anchor=north},legend columns=2]
\addplot[omDef,thick,mark=*,mark size=1.3pt] table[x=group,y=pred]{figdata/hft/06-adverse-selection-and-toxicity/calibration.csv};
\addplot[omThm,thick,only marks,mark=square*,mark size=1.5pt,error bars/.cd,y dir=both,y explicit]
  table[x=group,y=real,y error=se]{figdata/hft/06-adverse-selection-and-toxicity/calibration.csv};
\draw[gray] (axis cs:0.5,0) -- (axis cs:5.5,0);
\legend{mean forecast,realised mark-out ($\pm$ one standard error)}
\end{axis}
\end{tikzpicture}

\omcaption{Out of sample: the base quoter's fills on six test sessions, sorted into fifths by the fitted forecast of their ten-second mark-out;
mean forecast and realised mark-out per fifth. Data: \texttt{hf\_toxicity.calibration}.}\label{fig:hf:adverse-selection-and-toxicity:calibration}
\end{omfigure}

\section{Responses: widen, fade, skew, refuse}

A score is worth what the market maker does with it. There are four responses, and the market structure decides which are available.

\begin{definition}[Defensive widening]\label{def:hf:adverse-selection-and-toxicity:widening}
\emph{Defensive widening}\index{defensive widening} moves a quote away from the \omterm{def:hf:fair-value:fair}{fair price} on the side, or both sides, where adverse selection is
expected, keeping it in the book at a price that compensates for the expected mark-out; in a one-tick book it means stepping back a tick from the
best price.
\end{definition}

Fading withdraws the quote altogether; from the router's side, One Quant Book 10, chapter 18 calls the quote that disappears as an order arrives a
\emph{quote fade}. Skewing moves both quotes in the direction of the expected move (chapter 7). Refusing is the dealer's version: a dealer that
knows its counterparty prices it out (client tiering, One Quant Book 9, chapter 24) or rejects it under last look (chapter 9). On an anonymous
exchange a one-tick market maker has two of these: fade, or widen by stepping back
(\cref{fig:hf:adverse-selection-and-toxicity:responses}).

\begin{omfigure}
\begin{tikzpicture}[font=\small,x=1cm,y=0.5cm]
\foreach \dx/\name in {0/base,4/fade,8/widen} {
  \draw[gray!60] (\dx,0) -- (\dx+3,0);
  \node[font=\footnotesize,anchor=north] at (\dx+1.5,-1.3) {\name};
  \fill[omDef!35] (\dx+0.3,0) rectangle (\dx+0.9,3.0);
  \fill[omDef!35] (\dx+0.9,0) rectangle (\dx+1.4,1.4);
  \fill[omThm!35] (\dx+1.6,0) rectangle (\dx+2.1,4.0);
  \fill[omThm!35] (\dx+2.1,0) rectangle (\dx+2.7,2.2);
  \node[font=\scriptsize,anchor=north] at (\dx+0.6,-0.1) {99};
  \node[font=\scriptsize,anchor=north] at (\dx+1.15,-0.1) {100};
  \node[font=\scriptsize,anchor=north] at (\dx+1.85,-0.1) {101};
  \node[font=\scriptsize,anchor=north] at (\dx+2.4,-0.1) {102};
}
\draw[black,very thick] (1.15,1.4) -- (1.15,2.2);
\draw[black,very thick] (1.85,4.0) -- (1.85,4.8);
\draw[black,very thick] (9.85,4.0) -- (9.85,4.8);
\draw[black,very thick] (8.6,3.0) -- (8.6,3.8);
\draw[black,very thick] (5.85,4.0) -- (5.85,4.8);
\node[font=\scriptsize,align=center,anchor=west] at (11.3,2.6) {our lot:\\thick bar};
\end{tikzpicture}

\omcaption{The two responses a one-lot market maker has in a one-tick book when its bid is scored toxic (the book: bid queues in blue, ask queues
in red, prices in ticks; our lot as a thick bar on top of a queue). Base: our lot rests at the best bid, 100, and at the best ask, 101. Fade: the
bid is withdrawn. Widen: the bid steps back to 99, where it is almost never filled.}\label{fig:hf:adverse-selection-and-toxicity:responses}
\end{omfigure}

The quoter of \cref{lst:hf:adverse-selection-and-toxicity:quoter} applies either response whenever a side's score exceeds a threshold. On the six
test sessions, at a threshold of zero (a side is withdrawn whenever its fill is forecast to lose):

\begin{center}
\small
\begin{tabular}{@{}lrrrrr@{}}
\toprule
response & shares & informed share & ten-second mark-out & messages & P\&L (sd)\\
& a session & of fills & (ticks a share, se) & a session & (\$ a session)\\
\midrule
none & 19,167 & 29.4\% & $-0.005$ (0.096) & 696 & $-75.00$ (178.23)\\
fade & 6,550 & 18.3\% & 0.136 (0.199) & 1,012 & 37.75 (134.30)\\
widen & 6,683 & 20.0\% & 0.112 (0.222) & 1,513 & $-2.08$ (105.59)\\
\bottomrule
\end{tabular}
\end{center}

Both responses remove a third of the informed share of fills, which the simulator measures exactly; their mark-outs improve by about 0.1 to 0.15
ticks a share, but the standard errors are of the same size, and six sessions cannot rank fading against widening. Both cost two thirds of the
volume and send many more messages. Varying the threshold shows why:

\begin{center}
\small
\begin{tabular}{@{}lrrrr@{}}
\toprule
fade above a score of & shares & informed share & mark-out (se) & messages\\
\midrule
0 & 6,550 & 18.3\% & 0.136 (0.199) & 1,012\\
0.2 & 10,733 & 17.9\% & 0.214 (0.146) & 952\\
0.4 & 15,017 & 24.1\% & 0.186 (0.123) & 841\\
0.6 & 18,017 & 30.1\% & $-0.060$ (0.103) & 740\\
\bottomrule
\end{tabular}
\end{center}

At a threshold of zero the score of a side crosses the threshold back and forth, and every crossing cancels the order and sends it to the back of
its queue when it returns; chapter 5 showed that the back of a queue is where the adverse fills are. A threshold of 0.2 keeps the informed share as
low with 64\% more volume. From 0.4 the market maker fades too rarely and the informed share returns.

\section{Toxicity is not permanent}

The same quoter's fills in the news window of each calibration session, and before and after it:

\begin{center}
\small
\begin{tabular}{@{}lccc@{}}
\toprule
& before the news & during the news & after the news\\
\midrule
fills a minute & 6.97 & 17.89 & 6.69\\
informed share of fills & 34.7\% & 23.6\% & 32.3\%\\
ten-second mark-out (ticks a share) & 0.025 & 0.059 & 0.132\\
\bottomrule
\end{tabular}
\end{center}

In this market the news window brings three times the noise flow as well as a faster efficient price, so the fills during it are less toxic, not
more: a rule that widens on every news window would give up the best flow of the session. The point generalises: toxicity is a property of the flow
at a time, not of an instrument, a venue or a client for ever. A score is refitted on recent fills, its decay is measured as for any predictor (One
Quant Book 7, chapter 13), and a counterparty tiered as toxic is re-tested.

\section{Tutorial: scoring the next fill}\label{tut:hf:adverse-selection-and-toxicity:tut}

\begin{tutorial}
\textbf{Goal.} Measure who is on the other side of a market maker's fills, build a \omterm{def:hf:adverse-selection-and-toxicity:score}{toxicity score} from observable features, and test fading and
widening on it. \textbf{End state:} the three tables and \cref{fig:hf:adverse-selection-and-toxicity:classes,fig:hf:adverse-selection-and-toxicity:calibration}.
\begin{tutsteps}
\item \textbf{Features from the feed.} \texttt{firm.toxicity.FlowTracker} updates on every message (the harness now passes the message behind each
update, \texttt{ctx.last}).
\omcode{code/firm/toxicity/firm_toxicity.py}{58}{75}{Signed flow against each side, the other side's queue share, quiet time and spread.}\label{lst:hf:adverse-selection-and-toxicity:features}
\item \textbf{Calibrate.} \texttt{hf\_toxicity.fitted()} regresses the ten-second mark-outs of the base quoter's fills on the features at each
fill, on six calibration sessions.
\item \textbf{Respond.} \texttt{ToxicQuoter} fades or widens a side whose score exceeds the threshold.
\omcode{code/firm/toxicity/firm_toxicity.py}{115}{137}{Score each side on every update; fade or step back the toxic side.}\label{lst:hf:adverse-selection-and-toxicity:quoter}
\item \textbf{Evaluate} with \texttt{responses()} and \texttt{fade\_sweep()} on six other sessions; \texttt{by\_class()} and \texttt{by\_period()} use
the simulator's truth.
\end{tutsteps}
\textbf{What to change next.} Add hysteresis (fade above 0.2, return below 0.1) and measure the messages saved; score the whole book's toxicity
with VPIN and compare it with the fill-level score.
\end{tutorial}

\section{Build: the toxicity toolkit}\label{bld:hf:adverse-selection-and-toxicity:toxicity}

\begin{build}
\bfield{Purpose} Grouped mark-outs, streaming toxicity features, a fitted score and responses that any Quoter can wrap.
\bfield{Interface} \texttt{grouped(markouts, qty, groups)}; \texttt{FlowTracker(window).update(t, last, ext)}, \texttt{.features(side)};
\texttt{ToxicityModel.fit(X, y)}, \texttt{.predict}, \texttt{.score}; \texttt{ToxicQuoter(model, mode, threshold, size, limit, record)}. In
\texttt{firm.mmharness}: \texttt{ctx.last} (the message behind each update) and the \texttt{informed} entry of \texttt{Result.extra} (evaluation only).
\bfield{Rules} Features use only what the feed has shown at the time. A score is fitted on sessions different from those it is judged on. The
simulator's counterparty classes never reach the Quoter.
\bfield{Acceptance tests} \texttt{code/firm/toxicity/tests/}: features by hand on a three-message feed (window, quiet time, spread); the fit
recovers a planted linear mark-out; grouped means by hand. \texttt{code/firm/mmharness/tests/}: the Quoter sees add, cancel and execution messages;
the counterparty log has one entry per fill and both classes.
\bfield{Stretch} Hysteresis; a logistic score of the informed probability; per-venue scores when \texttt{firm.exchsim} runs several venues.
\end{build}

\begin{omsources}
\item D.~Easley, M.~López de Prado and M.~O'Hara, ``Flow toxicity and liquidity in a high-frequency world'', \emph{Review of Financial Studies}
25(5), 2012.
\item T.~G.~Andersen and O.~Bondarenko, ``VPIN and the flash crash'', \emph{Journal of Financial Markets} 17, 2014.
\item R.~Battalio, S.~A.~Corwin and R.~Jennings, ``Can brokers have it all? On the relation between make-take fees and limit order execution
quality'', \emph{Journal of Finance} 71(5), 2016.
\end{omsources}

\section{Exercises}

\begin{exercise}[$\star$]\label{exo:hf:adverse-selection-and-toxicity:1}
32\% of a quoter's fills are informed, with a mark-out of $-0.95$ ticks; the rest mark out at $+0.55$. What is the mark-out of all its fills?
\end{exercise}

\begin{exercise}[$\star$]\label{exo:hf:adverse-selection-and-toxicity:2}
With the fitted forecast, compute the bid's score when 2 lots were sold against it in the last two seconds, the ask holds 80\% of the two best
queues, the mid last moved half a second ago and the spread is one tick. Does the quoter fade at a threshold of 0? Of 0.2?
\end{exercise}

\begin{exercise}[$\star$]\label{exo:hf:adverse-selection-and-toxicity:3}
A rule earns 0.2 ticks a share on 10,000 shares a session with a one-cent tick. What is its ten-second edge a session?
\end{exercise}

\begin{exercise}[$\star\star$]\label{exo:hf:adverse-selection-and-toxicity:4}
Why do fading and widening both send more messages than the base quoter?
\end{exercise}

\begin{exercise}[$\star\star$]\label{exo:hf:adverse-selection-and-toxicity:5}
Why does the other side's share of the queues carry almost all of the forecast in this market?
\end{exercise}

\begin{exercise}[$\star\star$]\label{exo:hf:adverse-selection-and-toxicity:6}
From \cref{fig:hf:adverse-selection-and-toxicity:calibration}, what is the realised mark-out of the fifth of fills with the best forecast, and is
the difference with the worst fifth larger than its standard errors?
\end{exercise}

\begin{exercise}[$\star\star\star$]\label{exo:hf:adverse-selection-and-toxicity:7}
\emph{Coding.} Run the fading quoter with hysteresis, \texttt{hysteresis(0.2, 0.0)}: a side is faded when its score exceeds 0.2 and restored only
when it falls below 0. Compare its messages, volume and informed share with fading above 0.2 without hysteresis.
\end{exercise}

\begin{exercise}[$\star\star\star$]\label{exo:hf:adverse-selection-and-toxicity:8}
\emph{Find the flaw.} ``VPIN is at its highest of the year: widen every quote on every venue until it falls.''
\end{exercise}

\section{Problem: Who Is on the Other Side}

\begin{problem}[{Weekend problem --- who is on the other side}]\label{pb:hf:adverse-selection-and-toxicity:1}
A one-lot market maker in a one-tick stock wants to stop trading with the people who know more than it does, without stopping trading.

\medskip\noindent\textbf{Part I --- Measuring.}
\begin{enumerate}
\item What groups can a market maker mark out by on an anonymous exchange, and on a request-for-quote platform?
\item What did Battalio, Corwin and Jennings find about venues?
\item Give the mark-outs of informed and uninformed fills and the informed share.
\item Why is the total mark-out so close to zero?
\end{enumerate}

\medskip\noindent\textbf{Part II --- Predicting.}
\begin{enumerate}[resume]
\item What is VPIN, and what did Andersen and Bondarenko find?
\item Define a \omterm{def:hf:adverse-selection-and-toxicity:score}{toxicity score} and list the chapter's features.
\item Give the fitted forecast and its in-sample and out-of-sample correlations.
\item What does the calibration by fifths show?
\end{enumerate}

\medskip\noindent\textbf{Part III --- Responding.}
\begin{enumerate}[resume]
\item Define \omterm{def:hf:adverse-selection-and-toxicity:widening}{defensive widening} and distinguish it from fading, skewing and refusing.
\item Which responses does a one-tick market maker have on an anonymous exchange?
\item Give the three responses' shares, informed shares and mark-outs.
\item Why can six sessions not rank fading against widening?
\end{enumerate}

\medskip\noindent\textbf{Part IV --- The verdict.}
\begin{enumerate}[resume]
\item State the \emph{named result}: the change in the informed share of fills and in the ten-second mark-out when the quoter fades on its score at
a threshold of 0.2, and the volume given up.
\item Why is a threshold of zero worse than 0.2?
\item What happens above 0.4?
\item What did the news window show about toxicity over time?
\item How would you re-test a counterparty tiered as toxic?
\item How would the responses change on a small-tick asset?
\item What would a dealer do with the same score?
\item In one sentence: what is a \omterm{def:hf:adverse-selection-and-toxicity:score}{toxicity score} for?
\end{enumerate}
\end{problem}

\section{Interview questions}

\begin{interviewq}[$\star$ \iqroles{trader}]\label{iq:hf:adverse-selection-and-toxicity:1}
Your fills on one venue mark out much worse than on another. List three explanations and how you would tell them apart.
\end{interviewq}

\begin{interviewq}[$\star\star$ \iqroles{researcher}]\label{iq:hf:adverse-selection-and-toxicity:2}
You only observe the mark-outs of fills you got. How does that bias a toxicity model, and what can you do about it?
\end{interviewq}

\begin{interviewq}[$\star\star$ \iqroles{researcher}]\label{iq:hf:adverse-selection-and-toxicity:3}
A toxicity forecast correlates 0.18 with realised mark-outs. Is that useful? What else do you need to know?
\end{interviewq}

\begin{interviewq}[$\star\star$ \iqroles{developer}]\label{iq:hf:adverse-selection-and-toxicity:4}
A score crosses its threshold several times a second. What does that do to the quoter, and how do you fix it in code?
\end{interviewq}

\begin{interviewq}[$\star\star$ \iqroles{risk}]\label{iq:hf:adverse-selection-and-toxicity:5}
A market maker with quoting obligations wants to fade whenever its score is high. What constrains it?
\end{interviewq}

\begin{interviewq}[$\star\star\star$ \iqroles{researcher}]\label{iq:hf:adverse-selection-and-toxicity:6}
A fraction $\pi$ of fills are informed and lose $L$; the rest gain $G$. A score flags informed fills with probability $a$ and uninformed ones with
probability $b$. Fading every flagged fill, when does the expected mark-out per fill improve, and what happens to volume?
\end{interviewq}
